\documentclass[11pt]{article}
\usepackage[margin=1in]{geometry}
\usepackage[T1]{fontenc}
\usepackage[utf8]{inputenc}
\usepackage{lmodern}
\usepackage{hyperref}
\usepackage{xcolor}
\usepackage{fancyvrb}
\usepackage{listings}
\usepackage{CJKutf8}
\lstset{basicstyle=\ttfamily\small,columns=fullflexible,breaklines=true,upquote=true,frame=single,framerule=0.2pt,tabsize=2}

\title{\texttt{\detokenize{model/layers.py}} (Q\&A Documentation)}
\author{}
\date{}
\begin{document}
\begin{CJK*}{UTF8}{gbsn}
\maketitle

\paragraph{Q: What does this file implement at a high level?}
\textbf{A:} This file defines reusable neural building blocks (embeddings, FM interaction, MLP variants, cross/AFM layers) shared across multiple models.

\paragraph{Q: Why keep these layers in a separate file?}
\textbf{A:} It avoids duplication across models and keeps implementations consistent (e.g., embedding offsets, FM math, activation/dropout choices).

\paragraph{Q: What does this file export (public API)?}
\textbf{A:} The public API is the set of classes/functions other modules import (sometimes re-exported by package initializers). The key top-level definitions detected here are:

\begin{Verbatim}[fontsize=\small]
class FeaturesLinear  # methods: __init__, forward
class FeaturesEmbedding  # methods: __init__, forward
class SeqFeatureEmbedding  # methods: __init__, forward
class FactorizationMachine  # methods: __init__, forward
class MultiLayerPerceptron  # methods: __init__, forward
class DurationMultiLayerPerceptron  # methods: __init__, forward
class Swish  # methods: forward
class MultiLayerPerceptronD2Q  # methods: __init__, forward
class MultiLayerPerceptronTPM  # methods: __init__, forward
class AttentionalFactorizationMachine  # methods: __init__, forward
class CrossNetwork  # methods: __init__, forward
\end{Verbatim}

\paragraph{Q: What are the main dependencies (imports) and why do they matter?}
\textbf{A:} Imports show whether this file is model code (torch), data code (pandas), or evaluation/analysis (sklearn). They also determine runtime requirements.

\vspace{1mm}
\VerbatimInput[firstline=1,lastline=50,fontsize=\small]{/workspace/model/layers.py}


\paragraph{Q: What is \texttt{\detokenize{FeaturesLinear}} and what responsibility does it have?}
\textbf{A:} This class is a core unit in the pipeline. Its responsibility is inferred from its name and how run scripts call it.

\paragraph{Q: What are the important methods on \texttt{\detokenize{FeaturesLinear}}?}
\textbf{A:} In this repo, methods like forward/loss/predict define computation and training targets. The methods found are:

\begin{Verbatim}[fontsize=\small]
__init__
forward
\end{Verbatim}

\paragraph{Q: What does the implementation look like for \texttt{\detokenize{FeaturesLinear}}?}
\textbf{A:} Excerpt from the file:

\vspace{1mm}
\VerbatimInput[firstline=5,lastline=15,fontsize=\scriptsize]{/workspace/model/layers.py}


\paragraph{Q: Why is \texttt{\detokenize{FeaturesLinear}} needed by training scripts?}
\textbf{A:} Training scripts assume this class can be instantiated from the dataset description and can map a feature dictionary to outputs required by that script's loss. This separation makes it easy to swap models while reusing the same dataloader and metrics.

\paragraph{Q: What is \texttt{\detokenize{FeaturesEmbedding}} and what responsibility does it have?}
\textbf{A:} This class is a core unit in the pipeline. Its responsibility is inferred from its name and how run scripts call it.

\paragraph{Q: What are the important methods on \texttt{\detokenize{FeaturesEmbedding}}?}
\textbf{A:} In this repo, methods like forward/loss/predict define computation and training targets. The methods found are:

\begin{Verbatim}[fontsize=\small]
__init__
forward
\end{Verbatim}

\paragraph{Q: What does the implementation look like for \texttt{\detokenize{FeaturesEmbedding}}?}
\textbf{A:} Excerpt from the file:

\vspace{1mm}
\VerbatimInput[firstline=18,lastline=31,fontsize=\scriptsize]{/workspace/model/layers.py}


\paragraph{Q: Why is \texttt{\detokenize{FeaturesEmbedding}} needed by training scripts?}
\textbf{A:} Training scripts assume this class can be instantiated from the dataset description and can map a feature dictionary to outputs required by that script's loss. This separation makes it easy to swap models while reusing the same dataloader and metrics.

\paragraph{Q: What is \texttt{\detokenize{SeqFeatureEmbedding}} and what responsibility does it have?}
\textbf{A:} This class is a core unit in the pipeline. Its responsibility is inferred from its name and how run scripts call it.

\paragraph{Q: What are the important methods on \texttt{\detokenize{SeqFeatureEmbedding}}?}
\textbf{A:} In this repo, methods like forward/loss/predict define computation and training targets. The methods found are:

\begin{Verbatim}[fontsize=\small]
__init__
forward
\end{Verbatim}

\paragraph{Q: What does the implementation look like for \texttt{\detokenize{SeqFeatureEmbedding}}?}
\textbf{A:} Excerpt from the file:

\vspace{1mm}
\VerbatimInput[firstline=34,lastline=47,fontsize=\scriptsize]{/workspace/model/layers.py}


\paragraph{Q: Why is \texttt{\detokenize{SeqFeatureEmbedding}} needed by training scripts?}
\textbf{A:} Training scripts assume this class can be instantiated from the dataset description and can map a feature dictionary to outputs required by that script's loss. This separation makes it easy to swap models while reusing the same dataloader and metrics.

\paragraph{Q: What is \texttt{\detokenize{FactorizationMachine}} and what responsibility does it have?}
\textbf{A:} This class is a core unit in the pipeline. Its responsibility is inferred from its name and how run scripts call it.

\paragraph{Q: What are the important methods on \texttt{\detokenize{FactorizationMachine}}?}
\textbf{A:} In this repo, methods like forward/loss/predict define computation and training targets. The methods found are:

\begin{Verbatim}[fontsize=\small]
__init__
forward
\end{Verbatim}

\paragraph{Q: What does the implementation look like for \texttt{\detokenize{FactorizationMachine}}?}
\textbf{A:} Excerpt from the file:

\vspace{1mm}
\VerbatimInput[firstline=50,lastline=65,fontsize=\scriptsize]{/workspace/model/layers.py}


\paragraph{Q: Why is \texttt{\detokenize{FactorizationMachine}} needed by training scripts?}
\textbf{A:} Training scripts assume this class can be instantiated from the dataset description and can map a feature dictionary to outputs required by that script's loss. This separation makes it easy to swap models while reusing the same dataloader and metrics.

\paragraph{Q: What is \texttt{\detokenize{MultiLayerPerceptron}} and what responsibility does it have?}
\textbf{A:} This class is a core unit in the pipeline. Its responsibility is inferred from its name and how run scripts call it.

\paragraph{Q: What are the important methods on \texttt{\detokenize{MultiLayerPerceptron}}?}
\textbf{A:} In this repo, methods like forward/loss/predict define computation and training targets. The methods found are:

\begin{Verbatim}[fontsize=\small]
__init__
forward
\end{Verbatim}

\paragraph{Q: What does the implementation look like for \texttt{\detokenize{MultiLayerPerceptron}}?}
\textbf{A:} Excerpt from the file:

\vspace{1mm}
\VerbatimInput[firstline=68,lastline=87,fontsize=\scriptsize]{/workspace/model/layers.py}


\paragraph{Q: Why is \texttt{\detokenize{MultiLayerPerceptron}} needed by training scripts?}
\textbf{A:} Training scripts assume this class can be instantiated from the dataset description and can map a feature dictionary to outputs required by that script's loss. This separation makes it easy to swap models while reusing the same dataloader and metrics.

\paragraph{Q: What is \texttt{\detokenize{DurationMultiLayerPerceptron}} and what responsibility does it have?}
\textbf{A:} This class is a core unit in the pipeline. Its responsibility is inferred from its name and how run scripts call it.

\paragraph{Q: What are the important methods on \texttt{\detokenize{DurationMultiLayerPerceptron}}?}
\textbf{A:} In this repo, methods like forward/loss/predict define computation and training targets. The methods found are:

\begin{Verbatim}[fontsize=\small]
__init__
forward
\end{Verbatim}

\paragraph{Q: What does the implementation look like for \texttt{\detokenize{DurationMultiLayerPerceptron}}?}
\textbf{A:} Excerpt from the file:

\vspace{1mm}
\VerbatimInput[firstline=91,lastline=111,fontsize=\scriptsize]{/workspace/model/layers.py}


\paragraph{Q: Why is \texttt{\detokenize{DurationMultiLayerPerceptron}} needed by training scripts?}
\textbf{A:} Training scripts assume this class can be instantiated from the dataset description and can map a feature dictionary to outputs required by that script's loss. This separation makes it easy to swap models while reusing the same dataloader and metrics.

\paragraph{Q: What is \texttt{\detokenize{Swish}} and what responsibility does it have?}
\textbf{A:} This class is a core unit in the pipeline. Its responsibility is inferred from its name and how run scripts call it.

\paragraph{Q: What are the important methods on \texttt{\detokenize{Swish}}?}
\textbf{A:} In this repo, methods like forward/loss/predict define computation and training targets. The methods found are:

\begin{Verbatim}[fontsize=\small]
forward
\end{Verbatim}

\paragraph{Q: What does the implementation look like for \texttt{\detokenize{Swish}}?}
\textbf{A:} Excerpt from the file:

\vspace{1mm}
\VerbatimInput[firstline=114,lastline=116,fontsize=\scriptsize]{/workspace/model/layers.py}


\paragraph{Q: Why is \texttt{\detokenize{Swish}} needed by training scripts?}
\textbf{A:} Training scripts assume this class can be instantiated from the dataset description and can map a feature dictionary to outputs required by that script's loss. This separation makes it easy to swap models while reusing the same dataloader and metrics.

\paragraph{Q: What is \texttt{\detokenize{MultiLayerPerceptronD2Q}} and what responsibility does it have?}
\textbf{A:} This class is a core unit in the pipeline. Its responsibility is inferred from its name and how run scripts call it.

\paragraph{Q: What are the important methods on \texttt{\detokenize{MultiLayerPerceptronD2Q}}?}
\textbf{A:} In this repo, methods like forward/loss/predict define computation and training targets. The methods found are:

\begin{Verbatim}[fontsize=\small]
__init__
forward
\end{Verbatim}

\paragraph{Q: What does the implementation look like for \texttt{\detokenize{MultiLayerPerceptronD2Q}}?}
\textbf{A:} Excerpt from the file:

\vspace{1mm}
\VerbatimInput[firstline=118,lastline=137,fontsize=\scriptsize]{/workspace/model/layers.py}


\paragraph{Q: Why is \texttt{\detokenize{MultiLayerPerceptronD2Q}} needed by training scripts?}
\textbf{A:} Training scripts assume this class can be instantiated from the dataset description and can map a feature dictionary to outputs required by that script's loss. This separation makes it easy to swap models while reusing the same dataloader and metrics.

\paragraph{Q: What is \texttt{\detokenize{MultiLayerPerceptronTPM}} and what responsibility does it have?}
\textbf{A:} This class is a core unit in the pipeline. Its responsibility is inferred from its name and how run scripts call it.

\paragraph{Q: What are the important methods on \texttt{\detokenize{MultiLayerPerceptronTPM}}?}
\textbf{A:} In this repo, methods like forward/loss/predict define computation and training targets. The methods found are:

\begin{Verbatim}[fontsize=\small]
__init__
forward
\end{Verbatim}

\paragraph{Q: What does the implementation look like for \texttt{\detokenize{MultiLayerPerceptronTPM}}?}
\textbf{A:} Excerpt from the file:

\vspace{1mm}
\VerbatimInput[firstline=139,lastline=158,fontsize=\scriptsize]{/workspace/model/layers.py}


\paragraph{Q: Why is \texttt{\detokenize{MultiLayerPerceptronTPM}} needed by training scripts?}
\textbf{A:} Training scripts assume this class can be instantiated from the dataset description and can map a feature dictionary to outputs required by that script's loss. This separation makes it easy to swap models while reusing the same dataloader and metrics.

\paragraph{Q: What is \texttt{\detokenize{AttentionalFactorizationMachine}} and what responsibility does it have?}
\textbf{A:} This class is a core unit in the pipeline. Its responsibility is inferred from its name and how run scripts call it.

\paragraph{Q: What are the important methods on \texttt{\detokenize{AttentionalFactorizationMachine}}?}
\textbf{A:} In this repo, methods like forward/loss/predict define computation and training targets. The methods found are:

\begin{Verbatim}[fontsize=\small]
__init__
forward
\end{Verbatim}

\paragraph{Q: What does the implementation look like for \texttt{\detokenize{AttentionalFactorizationMachine}}?}
\textbf{A:} Excerpt from the file:

\vspace{1mm}
\VerbatimInput[firstline=160,lastline=185,fontsize=\scriptsize]{/workspace/model/layers.py}


\paragraph{Q: Why is \texttt{\detokenize{AttentionalFactorizationMachine}} needed by training scripts?}
\textbf{A:} Training scripts assume this class can be instantiated from the dataset description and can map a feature dictionary to outputs required by that script's loss. This separation makes it easy to swap models while reusing the same dataloader and metrics.

\paragraph{Q: What is \texttt{\detokenize{CrossNetwork}} and what responsibility does it have?}
\textbf{A:} This class is a core unit in the pipeline. Its responsibility is inferred from its name and how run scripts call it.

\paragraph{Q: What are the important methods on \texttt{\detokenize{CrossNetwork}}?}
\textbf{A:} In this repo, methods like forward/loss/predict define computation and training targets. The methods found are:

\begin{Verbatim}[fontsize=\small]
__init__
forward
\end{Verbatim}

\paragraph{Q: What does the implementation look like for \texttt{\detokenize{CrossNetwork}}?}
\textbf{A:} Excerpt from the file:

\vspace{1mm}
\VerbatimInput[firstline=187,lastline=207,fontsize=\scriptsize]{/workspace/model/layers.py}


\paragraph{Q: Why is \texttt{\detokenize{CrossNetwork}} needed by training scripts?}
\textbf{A:} Training scripts assume this class can be instantiated from the dataset description and can map a feature dictionary to outputs required by that script's loss. This separation makes it easy to swap models while reusing the same dataloader and metrics.

\paragraph{Q: Example: end-to-end data flow involving this file}
\textbf{A:} Core pipeline shape (schematic):

\begin{Verbatim}[fontsize=\small]
for features, label in dataloaders["train"]:
    out = model(features)
    loss = ...
    loss.backward(); optimizer.step()
\end{Verbatim}

\paragraph{Q: Full source code of the Python file}
\textbf{A:} The complete file is included verbatim for reference.

\VerbatimInput[fontsize=\scriptsize]{/workspace/model/layers.py}

\end{CJK*}
\end{document}